\section*{Hoofdstuk \ref{ch:g9:ions} --- Ionen en ionenoplossingen}

\begin{solution}{exo:g9:ions:1}
Een \omterm{def:g7:atoms-and-molecules:atom}{atoom} of een groep \omterm{def:g7:atoms-and-molecules:atom}{atomen} die \omterm{def:g9:inside-the-atom:nucleus}{elektronen} heeft afgestaan of
opgenomen en een lading draagt. Een \omterm{def:g9:ions:ion}{kation} is positief (\omterm{def:g9:inside-the-atom:nucleus}{elektronen}
afgestaan), een \omterm{def:g9:ions:ion}{anion} negatief (\omterm{def:g9:inside-the-atom:nucleus}{elektronen} opgenomen).
\end{solution}

\begin{solution}{exo:g9:ions:2}
\ce{Mg^{2+}}: 12 \omterm{def:g9:inside-the-atom:proton}{protonen}, 10 \omterm{def:g9:inside-the-atom:nucleus}{elektronen}.
\end{solution}

\begin{solution}{exo:g9:ions:3}
\ce{KCl}, \ce{MgCl2}, \ce{CaO}.
\end{solution}

\begin{solution}{exo:g9:ions:4}
Zout water bevat \omterm{def:g9:ions:ion}{ionen}, \ce{Na+(aq)} en \ce{Cl-(aq)}, die bewegen en
lading dragen. Suikermoleculen dragen geen lading.
\end{solution}

\begin{solution}{exo:g9:ions:5}
\ce{CuSO4}, \ce{AlCl3}, \ce{Na2CO3}.
\end{solution}

\begin{solution}{exo:g9:ions:6}
IJzer(III), \ce{Fe^{3+}}: \ce{Fe^{3+}(aq) + 3OH-(aq) -> Fe(OH)3(s)}.
\end{solution}

\begin{solution}{exo:g9:ions:7}
Sulfaat: 2 \omterm{def:g9:inside-the-atom:nucleus}{elektronen} te veel. Nitraat: 1.
\end{solution}

\begin{solution}{exo:g9:ions:8}
Wit met zilvernitraat: \ce{Cl-}. Blauw met natriumhydroxide:
\ce{Cu^{2+}}. Koper(II)chloride, \ce{CuCl2}.
\end{solution}

\begin{solution}{exo:g9:ions:9}
Vier: aan elke kant één (links, rechts, boven, onder).
\end{solution}

\begin{solution}{exo:g9:ions:10}
\ce{Fe^{2+}(aq) + 2OH-(aq) -> Fe(OH)2(s)}.
\end{solution}

\begin{solution}{exo:g9:ions:11}
Zinkchloride. Natriumionen geven met natriumhydroxide geen
\omterm{def:g9:ions:precipitate}{neerslag}; zinkionen geven een witte \omterm{def:g9:ions:precipitate}{neerslag}
volgens \ce{Zn^{2+}(aq) + 2OH-(aq) -> Zn(OH)2(s)}.
\end{solution}

\begin{solution}{exo:g9:ions:12}
\ce{CaCl2}: 2000 \omterm{def:g9:ions:ion}{chloride-ionen}. Ladingen: $1000 \times (+2) + 2000 \times
(-1) = 0$.
\end{solution}

\begin{solution}{pb:g9:ions:1}
\textbf{1.} IJzer(II), \ce{Fe^{2+}}: \ce{Fe(OH)2}.

\textbf{2.} IJzer(III), \ce{Fe^{3+}}: \ce{Fe(OH)3}.

\textbf{3.} \ce{Fe^{2+}} heeft er één meer: $26 - 2 = 24$ \omterm{def:g9:inside-the-atom:nucleus}{elektronen},
tegen $26 - 3 = 23$ voor \ce{Fe^{3+}}.

\textbf{4.} \ce{Fe^{3+}(aq) + 3OH-(aq) -> Fe(OH)3(s)}.

\textbf{5.} Het komt uit de put met ijzer(II)\omterm{def:g9:ions:ion}{ionen}; in contact met de
\omterm{def:g4:air-a-mixture-of-gases:air}{lucht} worden die ijzer(III)\omterm{def:g9:ions:ion}{ionen}, die de oranje vaste stof vormen.

\textbf{6.} \ce{Fe(OH)3}: $(+3) + 3 \times (-1) = 0$.

\textbf{7.} Niet als het uit de put komt: de ijzerionen zijn opgelost
en gaan door een filter heen. Zodra de oranje vaste stof gevormd is,
wel: het is een \omterm{def:g9:ions:precipitate}{neerslag}, die door een filter wordt tegengehouden.

\textbf{8.} Nee: een \omterm{def:g6:pure-substances-and-mixtures:mixture}{mengsel} (water en opgeloste \omterm{def:g9:ions:ion}{ionen}). Vers opgepompt
en helder is het homogeen.

\textbf{9.} $56 \times 1.67 \times 10^{-27} = \qty{9.35e-26}{kg}$.

\textbf{10.} $\qty{0.30}{mg} = \qty{0.30e-6}{kg} = \qty{3.0e-7}{kg}$.

\textbf{11.} Een \omterm{def:g9:ions:ion}{ion} verschilt van het \omterm{def:g7:atoms-and-molecules:atom}{atoom} door twee of drie
\omterm{def:g9:inside-the-atom:nucleus}{elektronen}, waarvan de massa verwaarloosbaar is naast de kern.

\textbf{12.} $3.0 \times 10^{-7} / 9.35 \times 10^{-26} \approx 3.2
\times 10^{18}$ ijzerionen per liter.
\end{solution}
